%% 第2巻 巻末（こたえ・記号カード・定数と単位・チェック表・つぎの巻へ・さくいん・奥付）（第7回編集委員会）
\makeatletter
%% ---------------- こたえ ----------------
\bkAnswerPages
\par\vspace{6mm}
\begin{tcolorbox}[bk base, colframe=bkTheme!60, colback=bkTheme!5, boxrule=0.8pt, arc=2mm, left=4mm, right=4mm, top=2.5mm, bottom=3mm]
{\sffamily\bfseries\color{bkDeep}\R{答}{こた}え\R{合}{あ}わせのしかた}\par\vspace{1.5mm}
\begin{enumerate}[label={\bkcircled{\arabic*}}, leftmargin=2.4\zw, itemsep=1pt, topsep=0pt]
\item まず\R{自分}{じ|ぶん}の\R{答}{こた}えを\R{全部}{ぜん|ぶ}\R{書}{か}いてから、ここを\R{見}{み}よう。
\item ちがっていたら、\R{答}{こた}えだけを\R{直}{なお}さずに、\R{同}{おな}じ\R{形}{かたち}の\R{例題}{れい|だい}にもどって、どの\R{行}{ぎょう}でちがったかをさがそう。
\item \R{単位}{たん|い}もいっしょに\R{確}{たし}かめよう。\R{数}{かず}が\R{合}{あ}っていても、\R{単位}{たん|い}がちがえば、ちがう\R{量}{りょう}です（2-1）。
\end{enumerate}
\end{tcolorbox}

%% ---------------- 記号カード ----------------
\bkKigoPages{%
\bkkigorow{$=$}{イコール}{\R{左}{ひだり}と\R{右}{みぎ}が\R{同}{おな}じ\R{大}{おお}きさ}{1-4}%
\bkkigorow{$-$}{マイナス}{ひき\R{算}{ざん}。0より\R{小}{ちい}さい\R{数}{かず}の\R{印}{しるし}にも\R{使}{つか}う（2-2）}{1-5}%
\bkkigorow{$\times$}{かける}{1つ\R{分}{ぶん}の\R{数}{かず}×いくつ\R{分}{ぶん}}{1-12}%
\bkkigorow{$\div$}{わる}{\R{同}{おな}じ\R{数}{かず}ずつ\R{分}{わ}ける。1あたりを\R{求}{もと}める}{1-14}%
\bkkigorow{$\dfrac{1}{2}$}{2\R{分}{ぶん}の1}{1を2つに\R{等}{ひと}しく\R{分}{わ}けた1つ\R{分}{ぶん}}{1-16}%
\bkkigorow{\sffamily cm$^2$}{\R{平方}{へい|ほう}センチメートル}{1\R{辺}{ぺん}1\,cmの\R{正方形}{せい|ほう|けい}の\R{広}{ひろ}さ}{1-21}%
\bkkigorow{$<$\ \ $>$\ \ $\leqq$\ \ $\geqq$}{\R{不等号}{ふ|とう|ごう}}{\R{大}{おお}きさをくらべる\R{記号}{き|ごう}。$\leqq$ は「\R{以下}{い|か}」、$\geqq$ は「\R{以上}{い|じょう}」}{2-2}%
}

%% 記号カードのつづき：手で書くときのちゅうい（この本は問題をノートに解く）
\def\bkbkhead#1{\par\addvspace{4mm}\noindent{\sffamily\bfseries\large\color{bkTheme}#1}\par\nobreak\vspace{2mm}}
\clearpage
\bkbkhead{\R{記号}{き|ごう}を\R{書}{か}くときのちゅうい}
{\small この\R{本}{ほん}では、\R{問題}{もん|だい}をノートに\R{解}{と}きます。\R{手}{て}で\R{書}{か}くと、ほかの\R{字}{じ}とまぎれやすい\R{記号}{き|ごう}があります。\R{書}{か}くときのコツです。}\par\vspace{2mm}
\begin{bkcenter}
\begin{tblr}{colspec={Q[c,m,wd=22mm] Q[l,m,wd=34mm] Q[l,m,wd=68mm]}, row{1}={bg=bkVol!14, font=\sffamily\bfseries\footnotesize}, hlines={0.4pt,bkLine}, rowsep=3pt, cells={font=\small}}
\R{記号}{き|ごう} & まぎれやすいもの & \R{書}{か}くときのコツ \\
$x$ & かけ\R{算}{ざん}の $\times$ & 2つの\R{曲}{ま}がった\R{線}{せん}を\R{背中}{せ|なか}あわせに\R{書}{か}く。\R{文字式}{も|じ|しき}では×を\R{書}{か}かない（2-12） \\
$1$ & \R{文字}{も|じ}の $l$（エル）、$I$（アイ） & 1は\R{上}{うえ}に\R{小}{ちい}さなはねをつける。$l$ は\R{下}{した}を\R{右}{みぎ}へ\R{曲}{ま}げる \\
$0$ & \R{文字}{も|じ}の O（オー） & 0はたて\R{長}{なが}に\R{書}{か}く \\
$2$ & \R{文字}{も|じ}の $z$ & $z$ はまん\R{中}{なか}に\R{横線}{よこ|せん}を1\R{本}{ぽん}\R{入}{い}れる \\
$t$ & たし\R{算}{ざん}の $+$ & $t$ は\R{下}{した}を\R{右}{みぎ}へ\R{曲}{ま}げる \\
$-$（マイナス） & \R{漢字}{かん|じ}の「\R{一}{いち}」、のばす\R{音}{おと}の「ー」 & \R{数}{かず}のすぐ\R{前}{まえ}に、\R{短}{みじか}く\R{書}{か}く。$-3$ のように、\R{数}{かず}とはなさない \\
$\div$ & $+$ & \R{横線}{よこ|せん}の\R{上下}{じょう|げ}に\R{点}{てん}を1つずつ \\
$\dfrac{2}{3}$ & 2/3 & \R{横線}{よこ|せん}は\R{分子}{ぶん|し}・\R{分母}{ぶん|ぼ}より\R{長}{なが}く。\R{分母}{ぶん|ぼ}は\R{下}{した} \\
$x^2$ & $x2$、$2x$ & 2は\R{右}{みぎ}かたに\R{小}{ちい}さく\R{書}{か}く（$x\times x$ のこと） \\
$\sqrt{2}$ & $\sqrt{\phantom{2}}2$ & \R{屋根}{や|ね}の\R{横線}{よこ|せん}を、\R{中}{なか}の\R{数}{かず}の\R{終}{お}わりまでのばす \\
$\leqq$ & $<$ & \R{下}{した}の\R{線}{せん}を2\R{本}{ほん}。「\R{以下}{い|か}」のときだけ\R{書}{か}く（2-2） \\
$5.0$ & $5$ & ずれのはばを\R{表}{あらわ}す0は\R{消}{け}さない（2-6） \\
\end{tblr}
\end{bkcenter}

%% ---------------- 定数・単位・公式 ----------------
\bkbacktitle{よく\R{使}{つか}う\R{定数}{てい|すう}・\R{単位}{たん|い}・\R{公式}{こう|しき}}
\bksubsection{1}{この\R{巻}{かん}で\R{使}{つか}う\R{単位}{たん|い}}
\begin{bkcenter}
\small
\begin{tblr}{colspec={Q[l,m,wd=22mm] Q[l,m,wd=30mm] Q[l,m,wd=57mm] Q[c,m,wd=12mm]}, row{1}={bg=bkVol!14, font=\sffamily\bfseries\footnotesize}, hlines={0.4pt,bkLine}, rowsep=1.4pt}
\R{量}{りょう} & \R{単位}{たん|い}（\R{読}{よ}み\R{方}{かた}） & ほかの\R{単位}{たん|い}との\R{関係}{かん|けい} & \R{章}{しょう} \\
\R{長}{なが}さ & m（メートル） & $\qty{1}{km}=\qty{1000}{m}$、$\qty{1}{m}=\qty{100}{cm}$、$\qty{1}{cm}=\qty{10}{mm}$ & 2-10 \\
\R{面積}{めん|せき} & m$^2$（\R{平方}{へい|ほう}メートル） & $\qty{1}{m^2}=\qty{10000}{cm^2}$ & 2-10 \\
\R{体積}{たい|せき} & m$^3$、L（リットル） & $\qty{1}{L}=\qty{1000}{cm^3}$、$\qty{1}{mL}=\qty{1}{cm^3}$、$\qty{1}{m^3}=\qty{1000}{L}$ & 2-7 \\
\R{質量}{しつ|りょう} & kg（キログラム） & $\qty{1}{kg}=\qty{1000}{g}$ & 2-10 \\
\R{時間}{じ|かん} & s（\R{秒}{びょう}） & 1\R{分}{ぷん}＝60\R{秒}{びょう}、\mbox{1\R{時間}{じ|かん}（h）＝3600\R{秒}{びょう}} & 2-10 \\
\R{速}{はや}さ & m/s、km/h & $\qty{1}{m/s}=\qty{3.6}{km/h}$。\R{時速}{じ|そく}を\R{秒速}{びょう|そく}にするには÷3.6 & {2-3\\2-10} \\
{\R{速}{はや}さの\\ふえ\R{方}{かた}} & m/s$^2$ & 1\R{秒}{びょう}ごとに、\R{速}{はや}さが\R{何}{なん}\,m/s ふえるか & 2-10 \\
\R{温度}{おん|ど} & ℃、K（ケルビン） & 0\,K は\R{約}{やく}$-273$\,℃。これより\R{低}{ひく}い\R{温度}{おん|ど}はない & 2-2 \\
\R{力}{ちから} & N（ニュートン） & \R{約}{やく}100\,gの\R{物}{もの}の\R{重}{おも}さが\R{約}{やく}1\,N & 2-15 \\
\R{電圧}{でん|あつ}・\R{電流}{でん|りゅう} & {V（ボルト）\\A（アンペア）} & \R{電圧}{でん|あつ}は\R{電気}{でん|き}をおしだす\R{強}{つよ}さ、\R{電流}{でん|りゅう}は\R{流}{なが}れる\R{量}{りょう} & 2-15 \\
\R{抵抗}{てい|こう} & Ω（オーム） & \R{電気}{でん|き}の\R{通}{とお}りにくさ。1\,Ω＝1\,V÷1\,A & 2-15 \\
\R{割合}{わり|あい} & \%（パーセント） & \R{割合}{わり|あい}の0.01が1\% & 2-5 \\
\R{情報}{じょう|ほう}の\R{量}{りょう} & ビット、バイト & 8ビット＝1バイト（256\R{通}{とお}りを\R{表}{あらわ}せる） & 2-17 \\
\end{tblr}
\end{bkcenter}

\bksubsection{2}{\R{世界}{せ|かい}\R{共通}{きょう|つう}の7つの\R{基本単位}{き|ほん|たん|い}（SI）}
\begin{bkcenter}
\small
\begin{tblr}{colspec={Q[l,m,wd=30mm] Q[c,m,wd=14mm] Q[l,m,wd=30mm] Q[l,m,wd=30mm]}, row{1}={bg=bkVol!14, font=\sffamily\bfseries\footnotesize}, hlines={0.4pt,bkLine}, rowsep=1.4pt}
\R{量}{りょう} & \R{記号}{き|ごう} & \R{読}{よ}み\R{方}{かた} & この\R{本}{ほん}で\R{出}{で}てくる\R{所}{ところ} \\
\R{長}{なが}さ & m & メートル & 2-1・2-10 \\
\R{質量}{しつ|りょう} & kg & キログラム & 2-10 \\
\R{時間}{じ|かん} & s & \R{秒}{びょう} & 2-10 \\
\R{電流}{でん|りゅう} & A & アンペア & 2-15 \\
\R{温度}{おん|ど} & K & ケルビン & 2-2 \\
\R{物質量}{ぶっ|しつ|りょう} & mol & モル & （\R{化学}{か|がく}で\R{使}{つか}う） \\
\R{光}{ひかり}の\R{強}{つよ}さ & cd & カンデラ & （\R{明}{あか}かりで\R{使}{つか}う） \\
\end{tblr}
\end{bkcenter}

\bksubsection{3}{\R{大}{おお}きな\R{数}{かず}・\R{小}{ちい}さな\R{数}{かず}を\R{表}{あらわ}す\R{記号}{き|ごう}（\R{接頭語}{せっ|とう|ご}）}
\begin{bkcenter}
\small
\begin{tblr}{colspec={Q[l,m,wd=15mm] *{7}{Q[c,m,wd=15.4mm]}}, column{1}={bg=bkVol!14, font=\sffamily\bfseries\footnotesize}, hlines={0.4pt,bkLine}, rowsep=1.4pt, colsep=2pt}
\R{記号}{き|ごう} & G & M & k & c & m & {\textmu} & n \\
\R{読}{よ}み\R{方}{かた} & ギガ & メガ & キロ & センチ & ミリ & マイクロ & ナノ \\
\R{倍}{ばい} & $10^9$ & $10^6$ & $10^3$ & $10^{-2}$ & $10^{-3}$ & $10^{-6}$ & $10^{-9}$ \\
\R{例}{れい} & GB & MW & {km\\kg} & cm & {mm\\mg} & {\textmu}g & nm \\
\end{tblr}
\end{bkcenter}

\clearpage
\bksubsection{4}{\R{計算}{けい|さん}のきまり}
\begin{bkcenter}
\small
\begin{tblr}{colspec={Q[l,m,wd=29mm] Q[l,m,wd=77mm] Q[c,m,wd=18mm]}, row{1}={bg=bkVol!14, font=\sffamily\bfseries\footnotesize}, hlines={0.4pt,bkLine}, rowsep=2pt}
きまり & \R{式}{しき}・\R{例}{れい} & \R{章}{しょう} \\
\R{計算}{けい|さん}の\R{順番}{じゅん|ばん} & かっこの\R{中}{なか}が\R{先}{さき}。×・÷は＋・－より\R{先}{さき}。\quad $2+3\times4=14$ & 2-11 \\
\R{入}{い}れかえ・\R{分配}{ぶん|ぱい} & $a+b=b+a$、\ $a\times b=b\times a$、\ $(a+b)\times c=a\times c+b\times c$ & 2-11・2-12 \\
まとめてわる & $(A+B)\div h=A\div h+B\div h$ & 2-12 \\
\R{正負}{せい|ふ}の\R{数}{かず} & $(-3)+(+5)=+2$、\ $(+2)-(+5)=-3$、\ $(-2)\times(-3)=+6$ & 2-2 \\
\R{分数}{ぶん|すう}のたし\R{算}{ざん} & \R{分母}{ぶん|ぼ}をそろえて（\R{通分}{つう|ぶん}）\R{分子}{ぶん|し}をたす。\ $\dfrac{1}{2}+\dfrac{1}{3}=\dfrac{3}{6}+\dfrac{2}{6}=\dfrac{5}{6}$ & 2-4 \\
\R{分数}{ぶん|すう}のわり\R{算}{ざん} & わる\R{数}{かず}の\R{分母}{ぶん|ぼ}と\R{分子}{ぶん|し}を\R{入}{い}れかえてかける。\ $\dfrac{2}{3}\div\dfrac{4}{5}=\dfrac{2}{3}\times\dfrac{5}{4}=\dfrac{5}{6}$ & 2-4 \\
\R{指数}{し|すう}のきまり & $x^m\times x^n=x^{m+n}$、\ $x^0=1$、\ $x^{-1}=\dfrac{1}{x}$、\ $x^{-2}=\dfrac{1}{x^2}$ & 2-7・2-12 \\
\R{平方根}{へい|ほう|こん} & $\sqrt{a}\times\sqrt{b}=\sqrt{ab}$。\ $\sqrt{2}\times\sqrt{3}=\sqrt{6}$ & 2-8 \\
\R{等式}{とう|しき}の\R{性質}{せい|しつ} & \R{両辺}{りょう|へん}に\R{同}{おな}じ\R{数}{かず}をたす・ひく・かける・わる（0でわるのはだめ）と、\R{等式}{とう|しき}のまま & 2-14 \\
\R{移項}{い|こう} & ＋3を\R{反対}{はん|たい}の\R{辺}{へん}へ\R{移}{うつ}すと－3。\ $x+3=8\ \to\ x=8-3$ & 2-14 \\
\end{tblr}
\end{bkcenter}

\bksubsection{5}{この\R{巻}{かん}の\R{公式}{こう|しき}}
\begin{bkcenter}
\small
\begin{tblr}{colspec={Q[l,m,wd=29mm] Q[l,m,wd=77mm] Q[c,m,wd=18mm]}, row{1}={bg=bkVol!14, font=\sffamily\bfseries\footnotesize}, hlines={0.4pt,bkLine}, rowsep=2pt}
\R{公式}{こう|しき} & \R{式}{しき} & \R{章}{しょう} \\
\R{速}{はや}さ & \R{速}{はや}さ＝\R{道}{みち}のり÷\R{時間}{じ|かん}、\ $v=\dfrac{x}{t}$、\ $x=vt$、\ $t=\dfrac{x}{v}$ & 2-3・2-14 \\
\R{変化量}{へん|か|りょう} & $\Delta x$＝（あと）$-$（まえ） & 2-3 \\
\R{割合}{わり|あい} & \R{割合}{わり|あい}＝くらべる\R{量}{りょう}÷もとにする\R{量}{りょう} & 2-5 \\
\R{平均}{へい|きん} & \R{平均}{へい|きん}＝\R{合計}{ごう|けい}÷\R{個数}{こ|すう} & 2-6 \\
\R{円周}{えん|しゅう} & \R{円周}{えん|しゅう}＝\R{直径}{ちょっ|けい}×$\pi$ & 2-9 \\
\R{比例}{ひ|れい}・\R{反比例}{はん|ぴ|れい} & $y=ax$、\ $y=\dfrac{a}{x}$（$xy$ が\R{一定}{いっ|てい}） & 2-15 \\
ばね & $F=kx$（$k$ はばね\R{定数}{てい|すう}、N/m） & 2-15 \\
オームの\R{法則}{ほう|そく} & $V=RI$、\ $I=\dfrac{V}{R}$ & 2-15 \\
\R{振}{ふ}り\R{子}{こ} & \R{糸}{いと}の\R{長}{なが}さが4\R{倍}{ばい}になると、\R{周期}{しゅう|き}は2\R{倍}{ばい}（\R{周期}{しゅう|き}は\R{長}{なが}さの $\sqrt{\phantom{2}}$ に\R{比例}{ひ|れい}） & 2-8 \\
\end{tblr}
\end{bkcenter}

\bksubsection{6}{よく\R{使}{つか}う\R{数}{かず}}
\begin{bkcenter}
\small
\begin{tblr}{colspec={Q[l,m,wd=44mm] Q[l,m,wd=44mm] Q[c,m,wd=14mm]}, row{1}={bg=bkVol!14, font=\sffamily\bfseries\footnotesize}, hlines={0.4pt,bkLine}, rowsep=1.4pt}
\R{名前}{な|まえ} & \R{値}{あたい} & \R{章}{しょう} \\
\R{円周率}{えん|しゅう|りつ} $\pi$ & 3.14159… & 2-9 \\
$\sqrt{2}$ & 1.41421… & 2-8 \\
$\sqrt{3}$ & 1.73205… & 2-8 \\
\R{光}{ひかり}の\R{速}{はや}さ & \R{約}{やく}$\qty{3.0e8}{m/s}$ & 2-7 \\
\R{落}{お}ちる\R{物}{もの}の\R{速}{はや}さのふえ\R{方}{かた} & \R{約}{やく}\qty{9.8}{m/s^2}（2-0 では\R{約}{やく}10） & {2-0\\2-10} \\
\end{tblr}
\end{bkcenter}

\clearpage
\bksubsection{7}{\R{表計算}{ひょう|けい|さん}とプログラムの\R{書}{か}き\R{方}{かた}}
\begin{bkcenter}
\small
\begin{tblr}{colspec={Q[l,m,wd=34mm] Q[l,m,wd=30mm] Q[l,m,wd=30mm] Q[l,m,wd=26mm]}, row{1}={bg=bkVol!14, font=\sffamily\bfseries\footnotesize}, hlines={0.4pt,bkLine}, rowsep=2pt}
\R{意味}{い|み} & \R{算数}{さん|すう}・\R{数学}{すう|がく} & \R{表計算}{ひょう|けい|さん}（2-18） & Python（2-19） \\
たす・ひく & $+$、$-$ & \texttt{+}、\texttt{-} & \texttt{+}、\texttt{-} \\
かける・わる & $\times$、$\div$ & \texttt{*}、\texttt{/} & \texttt{*}、\texttt{/} \\
\R{累乗}{るい|じょう} & $2^3$ & \texttt{=2\textasciicircum 3} & \texttt{2**3} \\
\R{式}{しき}の\R{始}{はじ}まり & ― & \texttt{=}\,で\R{始}{はじ}める & ― \\
\R{入}{い}れる（\R{代入}{だい|にゅう}） & ― & セルに\R{式}{しき}を\R{書}{か}く & \texttt{x = x + 1} \\
\R{等}{ひと}しいか & $=$ & \texttt{=} & \texttt{==} \\
\R{以上}{い|じょう} & $\geqq$ & \texttt{>=} & \texttt{>=} \\
ずらさない\R{番地}{ばん|ち} & ― & \texttt{\$B\$1} & ― \\
\end{tblr}
\end{bkcenter}
\begin{itemize}[leftmargin=1.4\zw, itemsep=2pt, topsep=4pt, label={\color{bkTheme}●}]
\item \R{表計算}{ひょう|けい|さん}で\R{式}{しき}をコピーすると、セルの\R{番地}{ばん|ち}もいっしょにずれます。ずらしたくない\R{番地}{ばん|ち}には \texttt{\$} をつけます（2-18）。
\item プログラムの \texttt{x = x + 1} は「$x$ に1をたして、$x$ に\R{入}{い}れなおす」という\R{命令}{めい|れい}で、\R{方程式}{ほう|てい|しき}ではありません（2-19）。
\item コンピュータは\R{数}{かず}を\R{有限}{ゆう|げん}のけたで\R{覚}{おぼ}えるので、0.1を10\R{回}{かい}たしても、ぴったり1にならないことがあります。\R{小数}{しょう|すう}の\R{答}{こた}えは、\R{少}{すこ}しずれていると\R{考}{かんが}えて\R{読}{よ}みます（2-19）。
\end{itemize}

%% ---------------- チェック表 ----------------
\bkbacktitle{できるようになったことチェック\R{表}{ひょう}}
{\small \R{各章}{かく|しょう}の「ふりかえり」を\R{集}{あつ}めた\R{表}{ひょう}です。\R{章}{しょう}を\R{学}{まな}び\R{終}{お}えたら1\R{回}{かい}め、\R{次}{つぎ}の\R{巻}{かん}に\R{進}{すす}む\R{前}{まえ}に2\R{回}{かい}めのチェックをしよう。できなかった\R{所}{ところ}は、ページにもどって\R{読}{よ}みなおします。}\par\vspace{2mm}
%% 自動生成（tools/mkchecklist.py）。各章の「ふりかえり」の項目を集めたもの
\begingroup\footnotesize
\def\bkcb{\tikz[baseline=-0.5ex]\draw[bkTheme, line width=0.7pt, rounded corners=0.5pt] (0,-0.3em) rectangle (0.8em,0.5em);}
\begin{longtblr}{colspec={Q[c,m,wd=9mm] X[l,m] Q[c,m,wd=10mm] Q[c,m,wd=9mm] Q[c,m,wd=9mm]}, width=\linewidth, rowhead=1, colsep=3pt,
  row{1}={bg=bkVol!14, font=\sffamily\bfseries\scriptsize}, hlines={0.4pt,bkLine}, rowsep=2.3pt}
\R{章}{しょう} & できるようになったこと & ページ & 1\R{回}{かい}め & 2\R{回}{かい}め \\
{\sffamily\bfseries\color{bkTheme}2-1} & 「かぞえる\R{数}{かず}」と「はかる\R{数}{かず}」のちがいを\R{説明}{せつ|めい}できる。 & \pageref{bk:2-1:kazoeru} & \bkcb & \bkcb \\
 & \R{量}{りょう}を「\R{数値}{すう|ち}×\R{単位}{たん|い}」で\R{表}{あらわ}し、\R{単位}{たん|い}をかえて\R{表}{あらわ}せる。 & \pageref{bk:2-1:tani} & \bkcb & \bkcb \\
 & \R{小数}{しょう|すう}の\R{位}{くらい}がわかり、10\R{倍}{ばい}・$\tfrac{1}{10}$ で\R{小数点}{しょう|すう|てん}を\R{動}{うご}かせる。 & \pageref{bk:2-1:kurai} & \bkcb & \bkcb \\
{\sffamily\bfseries\color{bkTheme}2-2} & 0より\R{小}{ちい}さい\R{数}{かず}を\R{負}{ふ}の\R{数}{かず}で\R{表}{あらわ}し、\R{数直線}{すう|ちょく|せん}にかける。 & \pageref{bk:2-2:fu} & \bkcb & \bkcb \\
 & \R{不等号}{ふ|とう|ごう}（$<$・$>$・$\leqq$・$\geqq$）で、\R{数}{かず}の\R{大}{おお}きさをくらべられる。 & \pageref{bk:2-2:futougou} & \bkcb & \bkcb \\
 & \R{正}{せい}の\R{数}{かず}・\R{負}{ふ}の\R{数}{かず}の ＋ － × ÷ の\R{計算}{けい|さん}ができる。 & \pageref{bk:2-2:keisan} & \bkcb & \bkcb \\
{\sffamily\bfseries\color{bkTheme}2-3} & \R{同}{おな}じ\R{種類}{しゅ|るい}の\R{量}{りょう}を、\R{単位}{たん|い}をそろえてたしたりひいたりできる。 & \pageref{bk:2-3:onaji} & \bkcb & \bkcb \\
 & ひき\R{算}{ざん}で\R{変化量}{へん|か|りょう}を\R{求}{もと}め、$\Delta$ を\R{使}{つか}って\R{書}{か}ける。 & \pageref{bk:2-3:henka} & \bkcb & \bkcb \\
 & わり\R{算}{ざん}で\R{速}{はや}さを、かけ\R{算}{ざん}で\R{道}{みち}のりを\R{求}{もと}められる。 & \pageref{bk:2-3:wari} & \bkcb & \bkcb \\
{\sffamily\bfseries\color{bkTheme}2-4} & わり\R{算}{ざん}の\R{答}{こた}えを\R{分数}{ぶん|すう}で\R{書}{か}き、\R{約分}{やく|ぶん}・\R{通分}{つう|ぶん}ができる。 & \pageref{bk:2-4:wari} & \bkcb & \bkcb \\
 & \R{分数}{ぶん|すう}の＋－×÷ができ、わり\R{算}{ざん}で\R{逆数}{ぎゃく|すう}をかけるわけを\R{説明}{せつ|めい}できる。 & \pageref{bk:2-4:naze} & \bkcb & \bkcb \\
 & 0と1のあいだの\R{数}{かず}をかけた・わった\R{答}{こた}えの\R{大}{おお}きさがわかる。 & \pageref{bk:2-4:chiisai} & \bkcb & \bkcb \\
{\sffamily\bfseries\color{bkTheme}2-5} & 2つの\R{量}{りょう}の\R{関係}{かん|けい}を\R{比}{ひ}で\R{表}{あらわ}し、\R{等}{ひと}しい\R{比}{ひ}をつくれる。 & \pageref{bk:2-5:hi} & \bkcb & \bkcb \\
 & \R{比}{ひ}をかけ\R{算}{ざん}でつないで、\R{歯車}{は|ぐるま}の\R{回転}{かい|てん}を\R{求}{もと}められる。 & \pageref{bk:2-5:tsunagu} & \bkcb & \bkcb \\
 & \R{割合}{わり|あい}を\R{百分率}{ひゃく|ぶん|りつ}で\R{表}{あらわ}し、「20\%\R{引}{び}き」の\R{量}{りょう}を\R{求}{もと}められる。 & \pageref{bk:2-5:biki} & \bkcb & \bkcb \\
{\sffamily\bfseries\color{bkTheme}2-6} & \R{最小}{さい|しょう}\R{目}{め}もりの10\R{分}{ぶん}の1まで\R{読}{よ}み、\R{平均}{へい|きん}を\R{求}{もと}められる。 & \pageref{bk:2-6:yomu} & \bkcb & \bkcb \\
 & \R{四捨五入}{し|しゃ|ご|にゅう}で、がい\R{数}{すう}をつくれる。 & \pageref{bk:2-6:gaisuu} & \bkcb & \bkcb \\
 & 「\qty{5.0}{cm}」と「\qty{5}{cm}」のちがいを $\pm$ で\R{説明}{せつ|めい}できる。 & \pageref{bk:2-6:yuukou} & \bkcb & \bkcb \\
{\sffamily\bfseries\color{bkTheme}2-7} & \R{同}{おな}じ\R{数}{かず}を\R{何回}{なん|かい}もかけた\R{数}{かず}を、\R{累乗}{るい|じょう}で\R{書}{か}いて\R{計算}{けい|さん}できる。 & \pageref{bk:2-7:ruijo} & \bkcb & \bkcb \\
 & \R{体積}{たい|せき}の\R{単位}{たん|い}や\R{大}{おお}きな\R{数}{かず}を、10の\R{累乗}{るい|じょう}で\R{表}{あらわ}せる。 & \pageref{bk:2-7:juu} & \bkcb & \bkcb \\
 & \R{指数}{し|すう}のきまり、0\R{乗}{じょう}、マイナスの\R{指数}{し|すう}がわかる。 & \pageref{bk:2-7:kimari} & \bkcb & \bkcb \\
{\sffamily\bfseries\color{bkTheme}2-8} & \R{面積}{めん|せき}から1\R{辺}{ぺん}を、2\R{乗}{じょう}ではさんで\R{求}{もと}められる。 & \pageref{bk:2-8:hasami} & \bkcb & \bkcb \\
 & $\sqrt{\phantom{2}}$ の\R{意味}{い|み}がわかり、\R{平方根}{へい|ほう|こん}を\R{求}{もと}められる。 & \pageref{bk:2-8:root} & \bkcb & \bkcb \\
 & \R{周期}{しゅう|き}が\R{長}{なが}さの $\sqrt{\phantom{2}}$ に\R{比例}{ひ|れい}することを\R{説明}{せつ|めい}できる。 & \pageref{bk:2-8:furiko} & \bkcb & \bkcb \\
{\sffamily\bfseries\color{bkTheme}2-9} & \R{分数}{ぶん|すう}を\R{小数}{しょう|すう}にすると、\R{終}{お}わるかくり\R{返}{かえ}すわけがわかる。 & \pageref{bk:2-9:kurikaesu} & \bkcb & \bkcb \\
 & \R{分数}{ぶん|すう}で\R{書}{か}けない\R{数}{かず}（\R{無理数}{む|り|すう}）があることがわかる。 & \pageref{bk:2-9:muri} & \bkcb & \bkcb \\
 & \R{円周}{えん|しゅう}＝\R{直径}{ちょっ|けい}×$\pi$ で、\R{円周}{えん|しゅう}や\R{直径}{ちょっ|けい}を\R{求}{もと}められる。 & \pageref{bk:2-9:pi} & \bkcb & \bkcb \\
{\sffamily\bfseries\color{bkTheme}2-10} & k・m・{\textmu} を\R{読}{よ}んで、\R{単位}{たん|い}を\R{書}{か}きかえられる。 & \pageref{bk:2-10:settou} & \bkcb & \bkcb \\
 & 「\R{単位}{たん|い}つきの1」で、\R{時間}{じ|かん}や\R{速}{はや}さを\R{換算}{かん|さん}できる。 & \pageref{bk:2-10:kansan} & \bkcb & \bkcb \\
 & \R{単位}{たん|い}を\R{計算}{けい|さん}して、\R{式}{しき}が\R{正}{ただ}しいかを\R{確}{たし}かめられる。 & \pageref{bk:2-10:keisan} & \bkcb & \bkcb \\
{\sffamily\bfseries\color{bkTheme}2-11} & ことばの\R{式}{しき}を、□や\R{文字}{も|じ}を\R{使}{つか}った\R{式}{しき}で\R{表}{あらわ}せる。 & \pageref{bk:2-11:moji} & \bkcb & \bkcb \\
 & \R{計算}{けい|さん}の\R{順番}{じゅん|ばん}のきまりと、そのわけがわかる。 & \pageref{bk:2-11:junban} & \bkcb & \bkcb \\
 & どんな\R{数}{かず}でも\R{成}{な}り\R{立}{た}つきまりを、\R{文字}{も|じ}で1\R{行}{ぎょう}に\R{書}{か}ける。 & \pageref{bk:2-11:kimari} & \bkcb & \bkcb \\
{\sffamily\bfseries\color{bkTheme}2-12} & ×をはぶき、÷を\R{分数}{ぶん|すう}にして、\R{文字式}{も|じ|しき}を\R{短}{みじか}く\R{書}{か}ける。 & \pageref{bk:2-12:yakusoku} & \bkcb & \bkcb \\
 & \R{同類項}{どう|るい|こう}をまとめたり、かっこをはずしたりできる。 & \pageref{bk:2-12:dourui} & \bkcb & \bkcb \\
 & \R{指数}{し|すう}のきまりを、\R{文字}{も|じ}を\R{使}{つか}って\R{書}{か}いて\R{使}{つか}える。 & \pageref{bk:2-12:shisu} & \bkcb & \bkcb \\
{\sffamily\bfseries\color{bkTheme}2-13} & \R{公式}{こう|しき}の\R{文字}{も|じ}に\R{数}{かず}を\R{代入}{だい|にゅう}して、\R{式}{しき}の\R{値}{あたい}を\R{求}{もと}められる。 & \pageref{bk:2-13:dainyu} & \bkcb & \bkcb \\
 & \R{単位}{たん|い}ごと\R{代入}{だい|にゅう}し、\R{単位}{たん|い}をそろえて\R{計算}{けい|さん}できる。 & \pageref{bk:2-13:tani} & \bkcb & \bkcb \\
 & マイナスの\R{数}{かず}・\R{分数}{ぶん|すう}・\R{式}{しき}を、かっこごと\R{代入}{だい|にゅう}できる。 & \pageref{bk:2-13:kakko} & \bkcb & \bkcb \\
{\sffamily\bfseries\color{bkTheme}2-14} & \R{等式}{とう|しき}の\R{性質}{せい|しつ}を、\R{天}{てん}びんで\R{説明}{せつ|めい}できる。 & \pageref{bk:2-14:tenbin} & \bkcb & \bkcb \\
 & $x=vt$ を $t$ や $v$ について\R{解}{と}ける。 & \pageref{bk:2-14:vt} & \bkcb & \bkcb \\
 & \R{移項}{い|こう}を\R{使}{つか}い、1\R{行}{ぎょう}ずつきまりを\R{書}{か}いて\R{変形}{へん|けい}できる。 & \pageref{bk:2-14:ikou} & \bkcb & \bkcb \\
{\sffamily\bfseries\color{bkTheme}2-15} & \R{比例}{ひ|れい}を\R{見}{み}つけ、$y=ax$ の\R{式}{しき}とグラフにできる。 & \pageref{bk:2-15:hirei} & \bkcb & \bkcb \\
 & ばねの $F=kx$、オームの\R{法則}{ほう|そく} $V=RI$ を\R{使}{つか}える。 & \pageref{bk:2-15:bane} & \bkcb & \bkcb \\
 & \R{反比例}{はん|ぴ|れい}を\R{見分}{み|わ}け、2\R{倍}{ばい}にするとどうなるか\R{答}{こた}えられる。 & \pageref{bk:2-15:hanpirei} & \bkcb & \bkcb \\
{\sffamily\bfseries\color{bkTheme}2-16} & \R{方程式}{ほう|てい|しき}をつくり、\R{表}{ひょう}でためして\R{解}{かい}をさがせる。 & \pageref{bk:2-16:houteishiki} & \bkcb & \bkcb \\
 & 1\R{次方程式}{じ|ほう|てい|しき}を1\R{行}{ぎょう}ずつ\R{解}{と}ける。 & \pageref{bk:2-16:tejun} & \bkcb & \bkcb \\
 & \R{表}{ひょう}・グラフ・\R{方程式}{ほう|てい|しき}で\R{解}{と}き、\R{確}{たし}かめられる。 & \pageref{bk:2-16:graph} & \bkcb & \bkcb \\
{\sffamily\bfseries\color{bkTheme}2-17} & 2\R{進数}{しん|すう}と10\R{進数}{しん|すう}を\R{行}{い}き\R{来}{き}できる。 & \pageref{bk:2-17:henkan} & \bkcb & \bkcb \\
 & ビットの\R{数}{かず}から、\R{何}{なん}\R{通}{とお}りの\R{数}{かず}を\R{表}{あらわ}せるかを\R{計算}{けい|さん}できる。 & \pageref{bk:2-17:bit} & \bkcb & \bkcb \\
 & 2\R{進数}{しん|すう}を4けたずつまとめて、16\R{進数}{しん|すう}で\R{書}{か}ける。 & \pageref{bk:2-17:juroku} & \bkcb & \bkcb \\
{\sffamily\bfseries\color{bkTheme}2-18} & 「＝」の\R{式}{しき}を\R{書}{か}き、コピーしてくり\R{返}{かえ}せる。 & \pageref{bk:2-18:shiki} & \bkcb & \bkcb \\
 & ずらしたくない\R{番地}{ばん|ち}に \$ をつけて\R{使}{つか}える。 & \pageref{bk:2-18:dollar} & \bkcb & \bkcb \\
 & \R{上}{うえ}の\R{行}{ぎょう}の\R{結果}{けっ|か}で\R{下}{した}の\R{行}{ぎょう}を\R{計算}{けい|さん}し、グラフにできる。 & \pageref{bk:2-18:kurikaesi} & \bkcb & \bkcb \\
{\sffamily\bfseries\color{bkTheme}2-19} & \R{変数}{へん|すう}と\R{代入}{だい|にゅう}のプログラムの\R{結果}{けっ|か}を\R{予想}{よ|そう}できる。 & \pageref{bk:2-19:dainyu} & \bkcb & \bkcb \\
 & くり\R{返}{かえ}しと「もし～なら」のプログラムを\R{読}{よ}める。 & \pageref{bk:2-19:kurikaesi} & \bkcb & \bkcb \\
 & \R{小数}{しょう|すう}のずれのわけと\R{対策}{たい|さく}を\R{説明}{せつ|めい}できる。 & \pageref{bk:2-19:zure} & \bkcb & \bkcb \\
\end{longtblr}
\endgroup


%% ---------------- つぎの巻へ（第3巻の予告と、この巻の道具がこの先どこで使われるか） ----------------
\bkbacktitle{つぎの\R{巻}{かん}へ}
\newcommand\bknextpill[1]{\bk@pill{#1}}
\def\bknextrow{1}
\newcommand\bknextcard[4]{%
  \begin{tcolorbox}[enhanced, nobeforeafter, box align=top, width=0.485\linewidth, colframe=bkVol3!75!black, colback=white, boxrule=0.9pt, arc=2.5mm,
    left=2.5mm, right=2.5mm, top=1.5mm, bottom=2mm, equal height group=bknext\bknextrow]
  \bknextpill{#1}\enspace{\sffamily\bfseries\small #2}\par\vspace{1mm}
  {\sffamily\bfseries\footnotesize\color{bkDeep3}#3}\par\vspace{1mm}
  {\footnotesize\setlength{\parindent}{0pt}{\sffamily\bfseries\color{bkInk}\R{使}{つか}う\R{道具}{どう|ぐ}}\enspace #4}
  \end{tcolorbox}}
\R{数}{かず}と\R{量}{りょう}のことばがそろいました。\R{第}{だい}3\R{巻}{かん}「\R{関数}{かん|すう}とデータ」では、この\R{巻}{かん}で\R{身}{み}につけた\R{文字}{も|じ}と\R{式}{しき}を\R{使}{つか}って、\R{表}{ひょう}・グラフ・\R{式}{しき}で\R{変化}{へん|か}のきまりをつかみます。\R{第}{だい}3\R{巻}{かん}の\R{章}{しょう}をのぞいてみましょう。\par\vspace{2.5mm}
\noindent\bknextcard{3-2}{\R{場所}{ば|しょ}を\R{数}{かず}で\R{表}{あらわ}す：\R{座標}{ざ|ひょう}}{「どこ」を2つの\R{数}{かず}で\R{伝}{つた}えよう}{マイナスの\R{数}{かず}と\R{数直線}{すう|ちょく|せん}（2-2）を、たてと\R{横}{よこ}の2\R{本}{ほん}にする。}\hfill
\bknextcard{3-3}{\R{関数}{かん|すう}は「\R{入}{い}れると\R{出}{で}てくる」ルール}{\R{表}{ひょう}・グラフ・\R{式}{しき}は、\R{同}{おな}じことを\R{言}{い}っている}{\R{文字}{も|じ}の\R{式}{しき}（2-11）と\R{代入}{だい|にゅう}（2-13）。}\par\vspace{1.2mm}
\def\bknextrow{2}\noindent\bknextcard{3-5}{グラフの\R{傾}{かたむ}きは「\R{変化}{へん|か}の\R{速}{はや}さ」}{\R{坂}{さか}が\R{急}{きゅう}なほど、\R{速}{はや}く\R{変}{か}わっている}{\R{変化量}{へん|か|りょう} $\Delta$（2-3）と\R{比例}{ひ|れい}（2-15）。}\hfill
\bknextcard{3-6}{2\R{本}{ほん}の\R{直線}{ちょく|せん}が\R{交}{まじ}わるところ：\R{連立方程式}{れん|りつ|ほう|てい|しき}}{ふたりが\R{出会}{で|あ}うのは、グラフの\R{交}{まじ}わる\R{点}{てん}}{\R{式}{しき}の\R{変形}{へん|けい}（2-14）と1\R{次方程式}{じ|ほう|てい|しき}（2-16）。}\par\vspace{1.2mm}
\def\bknextrow{3}\noindent\bknextcard{3-11}{\R{指数関数}{し|すう|かん|すう}：いまある\R{量}{りょう}に\R{比例}{ひ|れい}してふえる・へる}{\R{倍々}{ばい|ばい}でふえると、あっという\R{間}{ま}に\R{大}{おお}きくなる}{\R{累乗}{るい|じょう}と\R{指数}{し|すう}（2-7）と、\R{表計算}{ひょう|けい|さん}のくり\R{返}{かえ}し（2-18）。}\hfill
\bknextcard{3-13}{データのちらばり：\mbox{ヒストグラム}と\R{平均}{へい|きん}}{\R{何回}{なん|かい}もはかって、\R{本当}{ほん|とう}の\R{値}{あたい}に\R{近}{ちか}づこう}{はかった\R{数}{かず}の「ずれ」と\R{平均}{へい|きん}（2-6）。}\par
\bkfukidashi{delta}{この\R{巻}{かん}のことばは、この\R{先}{さき}の\R{全部}{ぜん|ぶ}の\R{巻}{かん}で\R{使}{つか}うよ。チェック\R{表}{ひょう}で\R{自信}{じ|しん}のない\R{章}{しょう}を\R{見}{み}なおしてから、\R{第}{だい}3\R{巻}{かん}へ\R{進}{すす}もう！}

\clearpage
\bkbkhead{この\R{巻}{かん}の\R{道具}{どう|ぐ}は、この\R{先}{さき}でこう\R{使}{つか}う}
{\small \R{第}{だい}2\R{巻}{かん}の\R{章}{しょう}と、その\R{道具}{どう|ぐ}を\R{使}{つか}う\R{第}{だい}3\R{巻}{かん}から\R{先}{さき}の\R{章}{しょう}です（\R{章}{しょう}のはじめの「\R{道}{みち}しるべ」の「この\R{章}{しょう}の\R{先}{さき}に」から\R{集}{あつ}めました）。}\par\vspace{3mm}
\begingroup\let\bkRorig\R\def\R#1#2{\mbox{\bkRorig{#1}{#2}}}% 第7回：表の中でルビの語を途中で切らない
%% 自動生成（tools/mknext.py）。各章の道しるべ「この章の先に」から、後ろの巻の章を集めたもの
\begin{longtblr}{colspec={Q[l,m,wd=58mm] X[l,m]}, width=\linewidth, rowhead=1, row{1}={bg=bkVol!14, font=\sffamily\bfseries\footnotesize},
  hlines={0.4pt,bkLine}, rowsep=1.5pt, cells={font=\footnotesize}}
この\R{巻}{かん}の\R{章}{しょう} & この\R{道具}{どう|ぐ}を\R{使}{つか}う\R{章}{しょう} \\
{\bknextpill{2-2}\ マイナスの\R{数}{かず}と\R{数直線}{すう|ちょく|せん}} & {\bknextpill{3-2}\ \R{場所}{ば|しょ}を\R{数}{かず}で\R{表}{あらわ}す：\R{座標}{ざ|ひょう}\\ \R{第}{だい}4\R{巻}{かん}から\R{先}{さき}：\bknextpill{4-10}\,\bknextpill{5-5}} \\
{\bknextpill{2-3}\ ＋ － × ÷ は、\R{何}{なに}を\R{表}{あらわ}しているの？} & {\bknextpill{3-5}\ グラフの\R{傾}{かたむ}きは「\R{変化}{へん|か}の\R{速}{はや}さ」\\ \R{第}{だい}4\R{巻}{かん}から\R{先}{さき}：\bknextpill{5-1}\,\bknextpill{5-9}} \\
{\bknextpill{2-5}\ \R{比}{ひ}・\R{割合}{わり|あい}・\R{百分率}{ひゃく|ぶん|りつ}} & {\R{第}{だい}4\R{巻}{かん}から\R{先}{さき}：\bknextpill{5-6}} \\
{\bknextpill{2-6}\ はかった\R{数}{かず}には「ずれ」がある} & {\bknextpill{3-14}\ ばらつきを1つの\R{数}{かず}で：\R{標準偏差}{ひょう|じゅん|へん|さ}\\ \R{第}{だい}4\R{巻}{かん}から\R{先}{さき}：\bknextpill{5-8}} \\
{\bknextpill{2-7}\ \R{何回}{なん|かい}もかける：\R{累乗}{るい|じょう}と\R{指数}{し|すう}} & {\R{第}{だい}4\R{巻}{かん}から\R{先}{さき}：\bknextpill{5-14}} \\
{\bknextpill{2-8}\ 2\R{回}{かい}かけると□になる\R{数}{かず}：\R{平方根}{へい|ほう|こん}} & {\bknextpill{3-10}\ ボールはいつその\R{高}{たか}さに\R{届}{とど}く？：2\R{次方程式}{じ|ほう|てい|しき}\\ \R{第}{だい}4\R{巻}{かん}から\R{先}{さき}：\bknextpill{4-4}\,\bknextpill{6-6}} \\
{\bknextpill{2-9}\ \R{分数}{ぶん|すう}で\R{書}{か}けない\R{数}{かず}がある？} & {\R{第}{だい}4\R{巻}{かん}から\R{先}{さき}：\bknextpill{4-2}\,\bknextpill{4-4}\,\bknextpill{5-14}} \\
{\bknextpill{2-10}\ \R{単位}{たん|い}の\R{世界}{せ|かい}ルール：SI\R{単位}{たん|い}} & {\R{第}{だい}4\R{巻}{かん}から\R{先}{さき}：\bknextpill{5-1}\,\bknextpill{5-5}} \\
{\bknextpill{2-11}\ なぜ\R{文字}{も|じ}を\R{使}{つか}うの？} & {\bknextpill{3-3}\ \R{関数}{かん|すう}は「\R{入}{い}れると\R{出}{で}てくる」ルール} \\
{\bknextpill{2-12}\ \R{文字式}{も|じ|しき}のやくそく} & {\bknextpill{3-8}\ かっこをはずす：\R{展開}{てん|かい}と\R{平方完成}{へい|ほう|かん|せい}\\ \R{第}{だい}4\R{巻}{かん}から\R{先}{さき}：\bknextpill{5-2}\,\bknextpill{5-4}} \\
{\bknextpill{2-13}\ \R{公式}{こう|しき}に\R{数}{かず}を\R{入}{い}れてみよう：\R{代入}{だい|にゅう}} & {\bknextpill{3-3}\ \R{関数}{かん|すう}は「\R{入}{い}れると\R{出}{で}てくる」ルール\\ \R{第}{だい}4\R{巻}{かん}から\R{先}{さき}：\bknextpill{5-2}\,\bknextpill{5-3}} \\
{\bknextpill{2-14}\ \R{式}{しき}を\R{変形}{へん|けい}して、ほしい\R{量}{りょう}をとり\R{出}{だ}そう} & {\bknextpill{3-6}\ 2\R{本}{ほん}の\R{直線}{ちょく|せん}が\R{交}{まじ}わるところ：\R{連立方程式}{れん|りつ|ほう|てい|しき}\\ \R{第}{だい}4\R{巻}{かん}から\R{先}{さき}：\bknextpill{5-0}} \\
{\bknextpill{2-15}\ \R{比例}{ひ|れい}と\R{反比例}{はん|ぴ|れい}} & {\bknextpill{3-3}\ \R{関数}{かん|すう}は「\R{入}{い}れると\R{出}{で}てくる」ルール\\ \bknextpill{3-4}\ グラフをかこう・\R{読}{よ}もう\\ \bknextpill{3-7}\ \R{反比例}{はん|ぴ|れい}のグラフ\\ \R{第}{だい}4\R{巻}{かん}から\R{先}{さき}：\bknextpill{5-12}\,\bknextpill{6-6}} \\
{\bknextpill{2-16}\ \R{何秒後}{なん|びょう|ご}に\R{追}{お}いつく？1\R{次方程式}{じ|ほう|てい|しき}} & {\bknextpill{3-4}\ グラフをかこう・\R{読}{よ}もう\\ \bknextpill{3-6}\ 2\R{本}{ほん}の\R{直線}{ちょく|せん}が\R{交}{まじ}わるところ：\R{連立方程式}{れん|りつ|ほう|てい|しき}\\ \bknextpill{3-10}\ ボールはいつその\R{高}{たか}さに\R{届}{とど}く？：2\R{次方程式}{じ|ほう|てい|しき}\\ \R{第}{だい}4\R{巻}{かん}から\R{先}{さき}：\bknextpill{6-2}} \\
{\bknextpill{2-17}\ コンピュータの\R{数}{かず}：2\R{進数}{しん|すう}と16\R{進数}{しん|すう}} & {\bknextpill{3-11}\ \R{指数関数}{し|すう|かん|すう}：いまある\R{量}{りょう}に\R{比例}{ひ|れい}してふえる・へる} \\
{\bknextpill{2-18}\ \R{表計算}{ひょう|けい|さん}で\R{計算}{けい|さん}しよう} & {\R{第}{だい}4\R{巻}{かん}から\R{先}{さき}：\bknextpill{5-8}\,\bknextpill{6-2}} \\
{\bknextpill{2-19}\ プログラムで\R{計算}{けい|さん}しよう} & {\R{第}{だい}4\R{巻}{かん}から\R{先}{さき}：\bknextpill{5-2}\,\bknextpill{6-2}} \\
\end{longtblr}

\endgroup

%% ---------------- さくいん ----------------
\bkbacktitle{さくいん}
{\small \R{数字}{すう|じ}はページです。その\R{語}{ご}がはじめて\R{出}{で}てくるページや、くわしく\R{説明}{せつ|めい}しているページをのせています。}\par\vspace{2mm}
\begin{multicols}{2}\small\setlength{\parskip}{0pt}
\directlua{bk.printidx()}
\end{multicols}

%% ---------------- 奥付 ----------------
\bkpadfour
\global\bk@decofalse\bkfolio{false}
\null\vfill
\noindent\begin{minipage}{\linewidth}\small
{\sffamily\bfseries\Large\color{bkDeep}\R{物理}{ぶつ|り}を\R{理解}{り|かい}するための\R{数学}{すう|がく}　\R{第}{だい}2\R{巻}{かん}}\par\vspace{1mm}
{\sffamily\bfseries\large \R{数}{かず}と\R{量}{りょう}のことば}\par\vspace{3mm}
{\color{bkVol}\hrule height 1pt}\vspace{3mm}
\begin{tabular}{@{}l@{\hspace{6mm}}l@{}}
\R{編者}{へん|しゃ} & 「\R{物理}{ぶつ|り}を\R{理解}{り|かい}するための\R{数学}{すう|がく}」\R{編集委員会}{へん|しゅう|い|いん|かい} \\
 & \R{委員長}{い|いん|ちょう}　\R{小学校長}{しょう|がっ|こう|ちょう}（まとめ\R{役}{やく}） \\
 & \R{委員}{い|いん}　\R{数学}{すう|がく}・\R{電気工学}{でん|き|こう|がく}・\R{機械工学}{き|かい|こう|がく}・\R{教育学}{きょう|いく|がく}の\R{大学教授}{だい|がく|きょう|じゅ}、\R{幼稚園教諭}{よう|ち|えん|きょう|ゆ}、\\
 & \phantom{\R{委員}{い|いん}　}\R{特別支援学校教諭}{とく|べつ|し|えん|がっ|こう|きょう|ゆ}、\R{心理士}{しん|り|し}、\R{科学解説者}{か|がく|かい|せつ|しゃ}、\R{小学校教諭}{しょう|がっ|こう|きょう|ゆ} \\
デザイン & ユニバーサルデザイン\R{担当委員}{たん|とう|い|いん} \\
\R{挿絵}{さし|え} & イラスト\R{担当委員}{たん|とう|い|いん} \\
\R{版}{はん} & \R{見本版}{み|ほん|ばん}（\R{第}{だい}7\R{回編集委員会}{かい|へん|しゅう|い|いん|かい}）　2026\R{年}{ねん}10\R{月}{がつ} \\
\end{tabular}\par\vspace{3mm}
{\color{bkVol}\hrule height 1pt}\vspace{2mm}
{\footnotesize この\R{本}{ほん}の\R{問題}{もん|だい}の\R{答}{こた}えは、\R{巻末}{かん|まつ}の「こたえ」にあります。くわしい\R{解}{と}き\R{方}{かた}は、\R{別冊}{べっ|さつ}の\R{問題集}{もん|だい|しゅう}と\R{指導書}{し|どう|しょ}にのせます。}
\end{minipage}
\makeatother
